\documentclass{article}
\usepackage[T1]{fontenc}
\usepackage{lmodern}
\usepackage{graphicx}
\usepackage{amsmath,amssymb}
\usepackage[a4paper,margin=2cm]{geometry}
\usepackage[absolute,overlay]{textpos}
\setlength{\TPHorizModule}{1mm}
\setlength{\TPVertModule}{1mm}
\usepackage[indonesian]{babel}
\usepackage{hyperref}
\setlength{\emergencystretch}{2em}
% unit: Halaman 33

\begin{document}

% === Halaman 33 (python/native) ===

33

• Contoh FSR adalah LFSR (Linear Feedback Shift Register) 

• Bit-bit di dalam register digeser satu bit ke kanan setiap kali pembangkitan bit 

luaran (= bit kunci alir atau keystream)

• Bit bn digeser ke tempat bn – 1, bn – 1 digeser ke bn – 2, dst, b2 digeser ke b1, b1 

terlempar ke luar

• Bit yang terlempar keluar (b1) menjadi bit luaran LFSR

• Selanjutnya dihitung bn yang baru, yaitu  bn  = bn  bn – 1  …  b2  b1

bn

bn - 1

...

b 4

b 3

b 2

b 1

Register Geser

Bit Keluaran

...

bn  = bn  bn – 1  …  b2  b1

b1

\end{document}
